\documentclass[12pt,a4paper]{report}

% ==========================================
% PACKAGES
% ==========================================
\usepackage[utf8]{inputenc}
\usepackage[T1]{fontenc}
\usepackage[francais]{babel}
\usepackage{geometry}
\usepackage{xcolor}
\usepackage{booktabs}
\usepackage{tabularx}
\usepackage{longtable}
\usepackage{array}
\usepackage{graphicx}
\usepackage{listings}
\usepackage{fancyhdr}
\usepackage{titlesec}
\usepackage{tcolorbox}
\usepackage{hyperref}
\usepackage{enumitem}
\usepackage{float}
\usepackage{tikz}
\usepackage[table]{xcolor}
\usepackage{caption}
\usetikzlibrary{shapes.geometric, arrows.meta, positioning}

\geometry{
  headheight=15pt,
  top=2.5cm, bottom=2.5cm,
  left=2.5cm, right=2.5cm
}

% ==========================================
% COULEURS GLOBALES
% ==========================================
\definecolor{primary}{RGB}{30,56,100}
\definecolor{secondary}{RGB}{46,117,182}
\definecolor{codegray}{RGB}{150,150,150}
\definecolor{accent}{RGB}{255,165,0}
\definecolor{lightblue}{RGB}{235,245,255}
\definecolor{successgreen}{RGB}{0,128,0}
\definecolor{captureOrange}{RGB}{210,120,0}
\definecolor{captureBg}{RGB}{255,252,235}

% ==========================================
% STYLE DES LISTINGS (CODE)
% ==========================================
\definecolor{lightgray}{rgb}{0.96,0.96,0.96}
\definecolor{darkgray}{rgb}{0.3,0.3,0.3}
\definecolor{codegreen}{RGB}{0,128,0}
\definecolor{codepurple}{RGB}{128,0,128}
\definecolor{codeblue}{RGB}{0,0,180}

% Style TypeScript / Node.js
\lstdefinestyle{tsstyle}{
  backgroundcolor=\color{lightgray},
  commentstyle=\color{codegreen}\itshape,
  keywordstyle=\color{codeblue}\bfseries,
  stringstyle=\color{codepurple},
  basicstyle=\ttfamily\small,
  breaklines=true,
  frame=single,
  rulecolor=\color{darkgray},
  numbers=left,
  numberstyle=\tiny\color{darkgray},
  numbersep=8pt,
  tabsize=2,
  showstringspaces=false,
  language=Java,
  morekeywords={async, await, const, let, interface, type, enum,
                export, import, from, require, process,
                describe, it, expect, beforeEach, afterEach},
  aboveskip=1em,
  belowskip=1em,
  captionpos=b
}

% ==========================================
% STYLE DES TITRES
% ==========================================
\titleformat{\chapter}[display]
  {\normalfont\huge\bfseries\color{primary}}
  {\chaptertitlename\ \thechapter}{20pt}{\Huge}
\titleformat{\section}
  {\normalfont\Large\bfseries\color{secondary}}
  {\thesection}{1em}{}
\titleformat{\subsection}
  {\normalfont\large\bfseries\color{primary}}
  {\thesubsection}{1em}{}

% ==========================================
% EN-TÊTE ET PIED DE PAGE
% ==========================================
\pagestyle{fancy}
\fancyhf{}
\fancyhead[L]{\small\color{secondary}Gestion de Flotte de Véhicules}
\fancyhead[R]{\small\color{secondary}Master 1 GIL --- 2025/2026}
\fancyfoot[C]{\thepage}
\renewcommand{\headrulewidth}{0.4pt}

% ==========================================
% BOÎTES INFO
% ==========================================
\tcbuselibrary{skins,breakable}
\newtcolorbox{infobox}{
  enhanced,
  colback=lightblue,
  colframe=secondary,
  fonttitle=\bfseries\color{white},
}

% ==========================================
% BOÎTE CAPTURE D'ÉCRAN
% ==========================================
\newcommand{\placeholder}[2]{%
  \begin{tikzpicture}
    \draw[dashed, gray!60, rounded corners=4pt]
      (0,0) rectangle (\textwidth-0.6cm, #1);
    \node[gray!60, font=\small\itshape] at (0.5*\textwidth-0.3cm, 0.5*#1)
      {[ Insérer capture : #2 ]};
  \end{tikzpicture}%
}

\newtcolorbox{capturebox}[1]{
  colback=captureBg, colframe=captureOrange,
  fonttitle=\bfseries\small\color{captureOrange},
  title={Capture d'écran -- #1},
  arc=4pt, boxrule=0.8pt,
  left=4pt, right=4pt, top=4pt, bottom=4pt,
}

% ==========================================
\begin{document}

% ==========================================
% PAGE DE TITRE
% ==========================================
\begin{titlepage}
\centering
\vspace*{2cm}

{\color{primary}\rule{\linewidth}{1pt}}
\vspace{0.5cm}

{\Huge\bfseries\color{primary} Rapport Technique \& Validation\par}
\vspace{0.4cm}
{\LARGE\color{secondary} Semaine 7 : Observabilité, CI/CD et E2E \par}
\vspace{0.3cm}

\vspace{0.5cm}
{\color{primary}\rule{\linewidth}{1pt}}

\vspace{2cm}

{\Large Système de Gestion de Flotte de Véhicules\par}
\vspace{1cm}
{\large Master 1 GIL --- 2025/2026\par}
\vspace{0.5cm}

\begin{tcolorbox}[colback=lightblue, colframe=secondary, width=0.9\textwidth]
    \centering
    \textbf{Informations du projet}\\[0.5em]
    \begin{tabular}{ll}
      \textbf{Auteurs :} & Kherbachi Rima \& Mouhoubi Leiticia\\
      \textbf{Formation :} & Master 1 GIL \\
      \textbf{Université :} & Université de Rouen Normandie \\
      \textbf{Année :} & 2025 -- 2026 \\
    \end{tabular}
  \end{tcolorbox}
\vspace{1cm}
{\large Avril 2026\par}

\vfill
\end{titlepage}

\newpage
\tableofcontents
\newpage

\chapter{Validation des Tests End-to-End (Cypress)}

Cette section présente les résultats des tests de bout en bout validant le front-end React.

\section{Détail des Suites et Explications Techniques}

\subsection{Smoke Test : \texttt{spec.cy.ts}}
\textbf{Objectif :} Valider l'intégrité globale de l'interface et les redirections de sécurité. \\
Ce test vérifie que l'accès à la racine de l'application redirige automatiquement vers le portail de connexion Keycloak (\texttt{/realms/fleet-management}) en l'absence de session valide. Il s'assure également que l'architecture de base (barre de navigation, structure DOM) charge sans plantage grave.

\subsection{Test 01 \& 02 : \texttt{auth.cy.ts} et \texttt{dashboard.cy.ts}}
\textbf{Objectif :} Valider la mécanique d'authentification et l'affichage des agrégats métier.
\begin{itemize}
    \item \textbf{Auth \& RBAC :} Vérification du flux d'accès selon le rôle (Admin, Manager, Utilisateur). Les tests simulent les tokens JWT injectés dans Keycloak.
    \item \textbf{Dashboard :} Contrôle de l'affichage cohérent des KPI (Total véhicules, requêtes API, alertes). Les tests s'assurent que les appels REST sous-jacents retournent un code 200 et que le React contextuel met bien à jour la vue.
\end{itemize}
\begin{center}
    \placeholder{6cm}{Cypress : Capture 01\_auth.cy.ts ou 02\_dashboard.cy.ts}
\end{center}

\subsection{Test 03 : \texttt{vehicules.cy.ts}}
\textbf{Objectif :} S'assurer du bon fonctionnement du registre d'immatriculation. \\
Ce script valide les opérations CRUD (Création, Lecture, Modification, Suppression). Le test mock la Gateway pour vérifier que le tableau HTML interprète correctement les données JSON retournées, notamment le basculement dynamique du statut d'actif à inactif via des boutons d'actions dédiés.
\begin{center}
    \placeholder{5cm}{Cypress : Capture 03\_vehicules.cy.ts}
\end{center}

\subsection{Test 04 : \texttt{conducteurs.cy.ts}}
\textbf{Objectif :} Vérfier la gestion des ressources humaines et le formatage des données. \\
Cypress teste ici le rendu du formulaire d'ajout d'un conducteur, vérifie que les erreurs de validation (champs vides) bloquent bien la soumission, et garantit que les licences de permis (A, B, C, D) extraites du tableau JSON sont correctement listées dans l'interface sous forme de "tags" visuels.
\begin{center}
    \placeholder{5cm}{Cypress : Capture 04\_conducteurs.cy.ts}
\end{center}

\subsection{Test 05 : \texttt{localisation.cy.ts}}
\textbf{Objectif :} Tester le système temps réel et la cartographie (Map). \\
C'est l'un des tests les plus complexes. Il intercepte les routes REST tout en gérant le cycle de vie du WebSocket. Il vérifie que :
\begin{itemize}
    \item Le conteneur Leaflet (\texttt{<div data-testid="leaflet-map">}) se charge sans erreur.
    \item Les WebSockets affichent l'indicateur visuel correctement (Connecté / Hors ligne).
    \item L'injection d'une position GPS factice (\texttt{veh-mock-123}) déclenche l'apparition d'un marqueur sur la carte et met à jour le panneau latéral avec la bonne vitesse et immatriculation.
\end{itemize}
\begin{center}
    \placeholder{6cm}{Cypress : Capture 05\_localisation.cy.ts}
\end{center}

\subsection{Test 06 : \texttt{maintenance.cy.ts}}
\textbf{Objectif :} Valider la gestion du cycle de vie des interventions mécaniques. \\
Le test interagit largement avec la logique métier : il vérifie le filtrage réactif (clic sur les filtres "PLANIFIEE" ou "TERMINEE"), le calcul correct des prix affichés (défilement de liste), et teste le bouton d'action "Démarrer" en s'assurant qu'une requête \texttt{PATCH} est bien émise au backend via l'API.
\begin{center}
    \placeholder{6cm}{Cypress : Capture 06\_maintenance.cy.ts}
\end{center}

\subsection{Test 07 : \texttt{rbac.cy.ts}}
\textbf{Objectif :} Garantir de manière transversale les comportements de navigation et d'interface partagés. \\
Ce test simule et vérifie les mécanismes de protection des routes (contrôle d'accès) et des composants globaux :
\begin{itemize}
    \item \textbf{Navigation protégée :} Vérification du panneau "Accès refusé" avec la redirection correcte vers le Dashboard via le bouton Retour.
    \item \textbf{Composants partagés (Modal) :} S'assure que les modals CRUD de l'application s'ouvrent, se ferment via un bouton, et se ferment avec la touche \texttt{Esc}.
    \item \textbf{Routes introuvables :} Redirection de sécurité automatique sur \texttt{/} lors de requêtes 404 (pages inexistantes par l'utilisateur).
\end{itemize}
\begin{center}
    \placeholder{5cm}{Cypress : Capture 07\_rbac.cy.ts}
\end{center}

\subsection{Test 08 : \texttt{status\_badge.cy.ts}}
\textbf{Objectif :} Vérifier la consistance de la logique d'Interface Utilisateur (Composants). \\
Ce test unitaire intégré cible spécifiquement la granularité visuelle. Il navigue ligne par ligne dans les tableaux React via \texttt{cy.within()} pour s'assurer que les administrateurs voient un menu déroulant (\texttt{select}) pour changer l'état d'un véhicule, là où les statuts non modifiables comme "Actif/Inactif" s'affichent sous forme de simples badges colorés (vert/rouge).
\begin{center}
    \placeholder{5cm}{Cypress : Capture 08\_status\_badge.cy.ts}
\end{center}

\subsection{Test 09 : \texttt{alertes.cy.ts}}
\textbf{Objectif :} Valider la supervision des incidents IoT et matériels. \\
Ce test intercepte l'API retournant le flux des alertes système. Il s'assure que :
\begin{itemize}
    \item L'indicateur numérique (badge) reflète exactement le nombre d'alertes non acquittées.
    \item Les modales de détails (bouton "Consulter") affichent bien le niveau de criticité et le message exact de l'incident.
    \item L'action "Acquitter" déclenche bien une requête de suppression (\texttt{DELETE}) et met à jour instantanément la liste sans rechargement de page.
\end{itemize}
\begin{center}
    \placeholder{6cm}{Cypress : Capture 09\_alertes.cy.ts}
\end{center}



\chapter{Observabilité et Performance}

La version de production a été outillée avec OpenTelemetry pour fournir une surveillance en temps réel de tous les composants de l'infrastructure.

\section{Grafana : Dashboard Métier et Système}
Affichage des compteurs globaux (CPU, RAM, requêtes) et des métriques métiers de la gestion de flotte. Le tableau de bord centralise la supervision en offrant une visibilité immédiate sur :
\begin{itemize}
    \item \textbf{Les métriques RED} (Rate, Errors, Duration) mesurant la charge et la fiabilité des API.
    \item \textbf{La consommation des ressources} Node.js et JVM pour prévenir les risques de fuites mémoire (OOM).
    \item \textbf{Les indicateurs métiers en temps réel}, comme le nombre de véhicules en cours de trajet, facilitant ainsi la prise de décision opérationnelle.
\end{itemize}
\begin{center}
    % \includegraphics[width=0.9\textwidth]{images/grafana_dashboard.png}
    \placeholder{8cm}{Dashboard Grafana (Métriques RED, Santé JVM/Node, Véhicules actifs)}
\end{center}

\section{Jaeger : Traçage Distribué}
Visualisation du suivi d'une requête traversant l'API Gateway vers les différents microservices. Grâce à l'instrumentation OpenTelemetry, Jaeger permet :
\begin{itemize}
    \item De reconstituer l'arbre d'exécution complet d'une requête HTTP (du frontend jusqu'à la base de données).
    \item D'identifier précisément les microservices responsables de goulots d'étranglement ou de latences réseau.
    \item De corréler les logs avec l'identifiant unique de la trace (TraceID).
\end{itemize}
\begin{center}
    \placeholder{7cm}{Interface Jaeger montrant une trace complète (Gateway -> Service -> DB)}
\end{center}

\section{Prometheus : Collecte des Métriques}
Vérification des targets enregistrés et de l'état de la collecte via OpenTelemetry Collector. Prometheus sert de cœur pour le monitoring :
\begin{itemize}
    \item \textbf{Scraping actif} des endpoints \texttt{/metrics} exposés par nos services.
    \item Stockage des données sous forme de séries temporelles (Time Series Database).
    \item Évaluation continue des alertes et maintien d'un historique robuste exploité ensuite par Grafana.
\end{itemize}
\begin{center}
    \placeholder{6cm}{Interface Prometheus (Targets Up/Down)}
\end{center}

\chapter{Échanges Événementiels et Tests de Charge}

\section{Apache Kafka : Monitoring des Topics}
Surveillance des messages échangés de manière asynchrone entre les microservices sur les principaux topics (ex: \texttt{fleet.events}). L'utilisation de Kafka garantit :
\begin{itemize}
    \item Le découplage fort entre le service de géolocalisation et les autres domaines métiers (maintenances, alertes).
    \item Une haute tolérance aux pannes et une persistance des événements.
    \item Une absorption efficace (buffering) des pics de sollicitations en cas d'envoi massif de coordonnées GPS.
\end{itemize}
\begin{center}
    \placeholder{6cm}{Console UI de Kafka (UI for Apache Kafka ou logs montrant les topics)}
\end{center}

\section{K6 : Script de Charge}
Exécution d'un script K6 simulant de multiples utilisateurs concurrents (Virtual Users). Cette étape de Load Testing permet :
\begin{itemize}
    \item De valider que le temps de réponse de la Gateway reste dans le seuil acceptable (P95 < 500ms).
    \item De simuler de véritables flux d'utilisation asynchrone (création de véhicules, positionnements).
    \item D'observer le comportement de l'infrastructure sous contrainte via l'observabilité.
\end{itemize}
\begin{center}
    \placeholder{7cm}{Terminal : Résultat final de l'exécution du script K6 (Summary report)}
\end{center}


\chapter{Déploiement Orchestré}

\section{Kubernetes \& Helm : Processus de Déploiement}
Le déploiement de l'ensemble de la flotte a été automatisé via la création d'un \textbf{Chart Helm} personnalisé. Ce chart regroupe tous les manifestes Kubernetes (Deployments, Services, ConfigMaps, Ingress) permettant de déployer l'infrastructure en une seule commande centralisée.

\subsection{Commandes de déploiement}

Voici les étapes exactes et le code exécuté pour déployer notre environnement sur le cluster Kubernetes (Minikube / K3s) :

\begin{lstlisting}[style=bashstyle, caption={Installation via Helm et vérification}]
# 1. Mise à jour des dépendances du Chart (si besoin)
helm dependency update ./fleet-chart

# 2. Déploiement de l'ensemble de l'architecture dans le namespace 'fleet'
helm upgrade --install fleet-release ./fleet-chart --namespace fleet --create-namespace

# 3. Vérification du statut de la release Helm
helm ls -n fleet

# 4. Suivi du démarrage et de l'état des Pods
kubectl get pods -n fleet --watch
\end{lstlisting}

\subsection{État des ressources orchestrées}

L'orchestration Kubernetes apporte la validation de la santé de nos conteneurs grâce aux sondes \texttt{liveness} et \texttt{readiness}. Une fois la commande \texttt{helm install} exécutée, nous pouvons observer la création structurée des déploiements (ReplicaSets) contrôlant nos pods.

\begin{center}
    \placeholder{6cm}{Terminal : Sortie de la commande \texttt{kubectl get pods -n fleet}}
\end{center}

\vspace{0.5cm}

\begin{center}
    \placeholder{6cm}{Terminal : Sortie de \texttt{kubectl get deployments -n fleet}}
\end{center}

\end{document}
